\documentclass{standalone}
\usepackage{circuitikz}
\usetikzlibrary{calc}
\definecolor{circuitnotemark}{HTML}{FE00FE}
\begin{document}
\begin{circuitikz}[american, line width=0.8pt]
\ctikzset{bipoles/length=1.2cm}
\coordinate (a1) at (0,0);
\coordinate (a3) at (2,0);
\draw (a1) to[R, l_=$R_{1}$, a^=$10\,\mathrm{k}\Omega$] (a3); % line 3
\path (3,-0.593) rectangle (4.27,0.593); % line 5
\node[anchor=west, inner sep=0, circuitnotemark, font=\large\bfseries] at (3,0) {X}; % line 5
\path (3,-6.249) rectangle (10.612,-0.407); % line 6
\node[anchor=west, inner sep=0, circuitnotemark, font=\large] at (3,-1) {X}; % line 6
\node[anchor=west, inner sep=0, circuitnotemark, font=\large] at (3,-1.423) {X}; % line 6
\node[anchor=west, inner sep=0, circuitnotemark, font=\large] at (3,-1.847) {X}; % line 6
\node[anchor=west, inner sep=0, circuitnotemark, font=\large] at (3,-2.27) {X}; % line 6
\node[anchor=west, inner sep=0, circuitnotemark, font=\large] at (3,-2.693) {X}; % line 6
\node[anchor=west, inner sep=0, circuitnotemark, font=\large] at (3,-3.117) {X}; % line 6
\node[anchor=west, inner sep=0, circuitnotemark, font=\large] at (3,-3.54) {X}; % line 6
\node[anchor=west, inner sep=0, circuitnotemark, font=\large] at (3,-3.963) {X}; % line 6
\node[anchor=west, inner sep=0, circuitnotemark, font=\large] at (3,-4.387) {X}; % line 6
\node[anchor=west, inner sep=0, circuitnotemark, font=\large] at (3,-4.81) {X}; % line 6
\node[anchor=west, inner sep=0, circuitnotemark, font=\large] at (3,-5.233) {X}; % line 6
\node[anchor=west, inner sep=0, circuitnotemark, font=\large] at (3,-5.657) {X}; % line 6
\path (9,-0.593) rectangle (10.27,0.593); % line 7
\node[anchor=west, inner sep=0, circuitnotemark, font=\large\bfseries] at (9,0) {X}; % line 7
\path (10,-8.112) rectangle (17.612,-0.407); % line 8
\node[anchor=west, inner sep=0, circuitnotemark, font=\large] at (10,-1) {X}; % line 8
\node[anchor=west, inner sep=0, circuitnotemark, font=\large] at (10,-1.593) {X}; % line 8
\node[anchor=west, inner sep=0, circuitnotemark, font=\large] at (10,-2.185) {X}; % line 8
\node[anchor=west, inner sep=0, circuitnotemark, font=\large] at (10,-2.778) {X}; % line 8
\node[anchor=west, inner sep=0, circuitnotemark, font=\large] at (10,-3.371) {X}; % line 8
\node[anchor=west, inner sep=0, circuitnotemark, font=\large] at (10,-3.963) {X}; % line 8
\node[anchor=west, inner sep=0, circuitnotemark, font=\large] at (10,-4.556) {X}; % line 8
\node[anchor=west, inner sep=0, circuitnotemark, font=\large] at (10,-5.149) {X}; % line 8
\node[anchor=west, inner sep=0, circuitnotemark, font=\large] at (10,-5.741) {X}; % line 8
\node[anchor=west, inner sep=0, circuitnotemark, font=\large] at (10,-6.334) {X}; % line 8
\node[anchor=west, inner sep=0, circuitnotemark, font=\large] at (10,-6.927) {X}; % line 8
\node[anchor=west, inner sep=0, circuitnotemark, font=\large] at (10,-7.519) {X}; % line 8
\node[anchor=south west, inner sep=2pt, circuitnotemark, font=\LARGE\bfseries] (circuittitle) at (current bounding box.north west) {X};
\path (circuittitle.south west) rectangle ++(4.298,0.853);
\end{circuitikz}
\end{document}
